\section{Systems and the two comparisons}\label{sec:systems}

\subsection{Our cascade}

Our system (\sweet~\cite{sweetsearch_repo}) stores three forms of each vector. Stage 1 uses a
512-bit sign code of the first 512 dimensions (L2-normalised before the sign), searched by Hamming
distance, either by a walk on an HNSW graph~\cite{malkov2020hnsw} built over the codes ($M = 64$,
$ef_{\text{construction}} = 800$, single-threaded build) or by an exact scan of all codes. The
production code scans up to 150,000 vectors and walks above that (a fixed cap, so that rankings do not
depend on the machine); in comparison (B) our cascade therefore scans on gcsn, d20k, AdvTest, FiQA and
dall, and walks on Quora and csn1m. In comparison (A) both generators run on every set up to the cap;
above it, only the walk exists. Stage 2
scores the stage-1 candidates with an int8 dot product over the same 512 dimensions. Stage 3 rescores
the best candidates with the full 768-dimensional vector stored as float16. The Hamming kernels use
AVX-512 \texttt{VPOPCNTDQ} where the CPU has it (Intel and AMD here), AVX2 otherwise, and NEON on ARM.
A short kernel-tuning phase ended before the dev selection and kept three changes: int8 rows mapped
from the index file (487~MB less anonymous memory at 1M, same speed), the code and int8 widths as build
parameters, and the AVX-512 Hamming path, which returns byte-identical result lists and cut dall scan
time on AMD from 287 to 279~\us{} and the 1M walk from 270 to 242~\us{} (dev, three interleaved rounds). Two
other candidate changes gave no gain and were dropped. Result lists on ARM are byte-identical to x86.
In comparison (B) the production entry point runs with $k = k_v$ results and one knob: the stage-1
candidate count $S_1$, with the walk's starting effort $ef = \max(S_1, 50)$. The grid includes values of
$S_1$ below the product's own floor for $k$, $\min(400, 1.5\max(40, 10k))$, which a bench-only setting
allows: \nm{ours-below-floor} of our \nm{ret-ours-n} frozen points at $\fr = 0.99$ use such a value, so they
are configurations the shipped product does not run. The shipped default ($S_1 = 400$) has no row of its
own. It runs in Node.js through a native addon; the timer
is inside the same process, around the call, which includes the product's orchestration and metadata
reads (\S\ref{sec:timing}). The walk adapts its effort per query: after the greedy descent it scales
$ef$ by 0.6 or 1.5 from the quality of the entry point and may stop early; $ef$ is the starting value.

\subsection{Rivals}

Table~\ref{tab:systems} lists the rival families, their libraries and the declared build grids.
All rivals are built from pinned sources with \texttt{-O3 -march=native} on each machine (FAISS at its
widest SIMD level for the CPU; hnswlib also with \texttt{-ffast-math}; DiskANN with its AVX2 kernels on
every x86 CPU): FAISS~\cite{douze2024faiss}
v1.15.1, USearch~\cite{vardanian2023usearch} v2.26.4, hnswlib~\cite{nmslib2017hnswlib} at a pinned commit,
DiskANN~\cite{subramanya2019diskann} at a pinned commit, Intel SVS~\cite{aguerrebere2023lvq,tepper2024leanvec}
v0.4.0 (LVQ and LeanVec run on Intel only), and ScaNN~\cite{guo2020scann,sun2023soar} 1.4.2 (Python
wheel; on ARM in a CPython 3.9 environment, the only aarch64 build). DiskANN has no ARM build. ScaNN
is timed through its Python binding. The binding and wrapper add a fixed cost per query, which we
measured with a one-vector index (\texttt{py\_overhead.py}, 200,000 calls, pinned): \nm{scann-ovh-intel}~\us{}
on Intel, \nm{scann-ovh-amd}~\us{} on AMD and \nm{scann-ovh-arm}~\us{} on Graviton (median). We do not
subtract it.

\begin{table}[t]\centering\small
\caption{Rival families in comparison (B), with the build grid of each set. $n$ = vectors; IVF lists
$\sqrt n$ and $4\sqrt n$. HNSW builds: $M \in \{16, 32, 64\}$, $ef_{\text{construction}} \in \{200, 800\}$.}\label{tab:systems}
\begin{tabularx}{\textwidth}{llX}\toprule
System & Library & Index and search\\\midrule
faiss\_exact & FAISS & exact float32 inner product (the latency reference)\\
faiss\_bin\_flat & FAISS & exact Hamming scan of 512-bit codes, float rescoring of a refine pool\\
faiss\_bin\_hnsw & FAISS & HNSW on 512-bit codes, refine pool rescored in float\\
faiss\_hnsw\_flat & FAISS & HNSW on float32\\
faiss\_hnsw\_sq8 & FAISS & HNSW on 8-bit scalar-quantised vectors, float refine\\
faiss\_ivfpq\_fs & FAISS & IVF with 4-bit PQ fast scan, float refine\\
faiss\_rabitq\_flat / \_ivf / \_hnsw & FAISS & RaBitQ~\cite{gao2024rabitq} codes, flat, IVF or HNSW, float refine\\
hnswlib\_f32 & hnswlib & HNSW on float32\\
usearch\_f16 / \_i8 & USearch & HNSW on float16 or int8 vectors\\
usearch\_b1\_rescore & USearch & HNSW on 1-bit codes, float rescoring of a refine pool\\
diskann\_vamana & DiskANN & in-memory Vamana graph, $R \in \{32, 64, 96\}$, $L \in \{100, 200\}$\\
svs\_vamana\_lvq / \_leanvec & SVS & Vamana with LVQ or LeanVec compression (Intel only)\\
scann & ScaNN & brute force ($n < 20$k), asymmetric hashing ($n < 100$k), tree + AH with and without SOAR\\\bottomrule
\end{tabularx}
\end{table}

\subsection{Comparison (A): a common rescoring procedure}\label{sec:compA}

Comparison (A) asks which candidate generator works best when everything after generation is the
same. Each generator returns its top $C$ vector ids. One shared C++ path then (1) scores all $C$
with the int8 dot product and keeps the best $H$ (stable sort; ties keep the generator order), and
(2) rescores those $H$ with the full float16 vectors in float64 and orders them. $H = C$ is the arm
without the int8 screen. The budget grid is declared: $C \in \{100, 200, 400, 1000\}$,
$H \in \{50, 100\}$ with $C \ge H$, plus $H = C$. Every generator runs at two efforts per budget: the
smallest effort that returns $C$ candidates ($ef = C$ for graphs) or, for partitions,
$\lceil 3\,C\,n_{\text{list}}/n\rceil$ lists (about three times $C$ vectors probed), and twice that. The
partition efforts are small: a partition generator returns $C$ candidates but not necessarily good
ones, so its (A) quality reflects this effort rule, not the best quality of the method. For our walk, $ef$ is the starting value of its adaptive rule. Nothing is selected per generator
in (A); we report every effort at every budget. The rival graphs of comparison (A) are built once per
set and model at the middle of their build grids ($M = 32$, $ef_{\text{construction}} = 800$ for HNSW
graphs; $R = 64$, $L = 200$ for DiskANN; $\sqrt n$ lists for IVF); ours uses its production build
($M = 64$, $ef_{\text{construction}} = 800$).

\paragraph{What is common and what is not.} The screen and the rescoring are common in their
\emph{output}: our two generators run the same steps through our addon kernels, and the shared C++ path
gives bit-identical retained lists to the addon on every dev set, on x86 and ARM. They are not common in
\emph{speed}. The rivals' path uses a scalar int8 dot product and a software float16 decoder; our addon
uses SIMD kernels with hardware float16 conversion. From the change in time between $H = 50$ and
$H = 1000$ at $C = 1000$, where the candidates are fixed, one rescored vector costs a median
\nm{rescore-ours-intel}~\us{} in ours and \nm{rescore-riv-intel}~\us{} in the shared C++ path on Intel
(\nm{rescore-x-intel}$\times$), \nm{rescore-x-amd}$\times$ on AMD and \nm{rescore-x-arm}$\times$ on Graviton. In
addition, our walk uses a denser graph ($M = 64$ against 32), adapts its effort per query, and is the only
process that asks for transparent huge pages. Comparison (A) therefore measures each generator
\emph{with its own implementation of the shared steps}; \S\ref{sec:results} gives the ratios with the
rescoring difference removed by this estimate. Generators: our walk, our
scan, FAISS binary HNSW and binary flat, USearch b1, hnswlib and FAISS HNSW on float32, FAISS RaBitQ
(flat and IVF), ScaNN and DiskANN.

\subsection{Comparison (B): as shipped}

In comparison (B) each library uses its own search and returns its top $k_v$ vectors. Where a
configuration refines (FAISS binary and quantised indexes, USearch b1), the refine step is one routine
of our harness for all of them: inner products with the float32 originals and a partial sort, not each
library's own refine class. The rival timer also includes one conversion of the query into the binary,
int8 and float64 forms that comparison (A) needs, also for systems that do not use them; this adds a
small fixed cost to every rival query, which we did not measure separately. $k_v$ is fixed per set and model on dev: the smallest number of vectors at which the exact
vector list holds 10 distinct documents for at least 99.9\% of dev queries (Quora systems return
$k_v + 1$ for the self-match). It ranges from 10 (BEIR) to 143 (gcsn with bge-base-en-v1.5). Each family's dev grid sweeps its search effort
(and refine depth where it has one); the operating point at each fraction is the fastest point that
meets it on dev (\S\ref{sec:protocol}). Comparison (B) answers the practical question: tuned on the
same dev queries, how fast is each system at the same retained quality?

\subsection{Memory, load and cold start}

Every timed job records the resident set size after loading and after the timed pass (on Linux from
\texttt{/proc}, on macOS from the task's physical footprint as well). A separate \emph{cold} job per
frozen point starts a new process after the page cache is dropped (\texttt{drop\_caches} on Linux,
\texttt{purge} on macOS) and measures index load time and the first query. Memory is resident memory after the timed pass: on Linux \texttt{Rss} (anonymous and file-backed pages, so
memory-mapped index files count when touched). Rivals that refine with the float32 originals hold them in
the harness ($n \times 768 \times 4$ bytes), as their library would. We do not report memory on the Mac:
the rival harness reads the float32 originals through temporary buffers, and macOS keeps those freed pages
in the process footprint, which inflates the rival numbers there. The same read path makes the cold
start of the refining rivals include a buffered copy of the originals, which a production server could
avoid by mapping them; we report cold start as measured and treat it as indicative. Build cost is reported for the 1M set from the builds that produced the timed indexes;
those builds ran concurrently under a memory-admission scheduler, so their wall times are indicative
only.
